% ==============================================================================
% CAPÍTULO 13: PITCH (Orquestador del Capítulo)
% ==============================================================================
\preportadacapitulo{13}{Pitch}

\newpage
\section{Capítulo 13: Pitch}\label{cap:pitch}

% [INSTRUCCIÓN]: Pitch para presentar la idea del proyecto. Contenido propio de esta versión; solo usa cifras ya citadas en el Capítulo 1 y no inventa precios ni cifras de mercado.
% [VERIFICAR]: antes de presentar, actualizar el estado del producto (frase marcada en 13.2.5) con los CU implementados a esa fecha (Capítulo~\hyperref[cap:producto]{12}) y ajustar el tiempo al que fije el docente o el jurado.

Este capítulo presenta el proyecto en formato de \textit{pitch}: una exposición breve y persuasiva que explica qué problema resuelve el sistema, cómo lo resuelve y por qué vale la pena construirlo. Incluye una versión corta (de 30 a 60 segundos), una versión completa con guion y tiempos, y las respuestas previstas a las preguntas más probables. Todas las cifras provienen de los datos de la DUBSS citados en la sección~1.2.

\subsection{13.1. Pitch de Ascensor}

El pitch de ascensor resume la idea en una sola intervención, de unos 45 segundos (alrededor de 110 palabras), para quien no conoce el proyecto.

\begin{quote}
Cada año, unos 20\,000 estudiantes de la UAGRM piden la Beca Comedor, pero solo hay cerca de 900 cupos y la DUBSS alcanza a revisar apenas 8\,000 carpetas de papel. Las otras 12\,000 ni siquiera se abren: quién obtiene la beca depende del orden de llegada o del azar, no de la necesidad.

Nosotros, en NEXWARE, lo cambiamos con un sistema web y móvil: el estudiante postula desde su celular, una inteligencia artificial calcula un puntaje de vulnerabilidad para las 20\,000 solicitudes, y la DUBSS audita solo a los finalistas mejor puntuados. Así, todos son evaluados con el mismo criterio y la beca llega a quien más la necesita.
\end{quote}

\subsection{13.2. Pitch Completo}

La versión completa dura unos 4 minutos y sigue el orden clásico de un pitch: problema, solución, cómo funciona, valor, avance y cierre con una petición. La Tabla~\ref{tab:pitch_tiempos} resume la secuencia y el tiempo de cada parte; el guion de cada una se desarrolla a continuación.

\begin{table}[H]
\centering
\caption{Estructura y tiempos del pitch completo}\label{tab:pitch_tiempos}
\begin{threeparttable}
\begin{tabular}{@{}l >{\raggedright\arraybackslash}p{3.0cm} >{\raggedright\arraybackslash}p{7.2cm} r@{}}
\toprule
Parte & Título & Mensaje central & Tiempo \\
\midrule
13.2.1 & Gancho & Una pregunta que muestra la injusticia del proceso actual. & 0:20 \\
13.2.2 & El problema & 20\,000 solicitudes, 900 cupos y 12\,000 carpetas sin abrir. & 0:40 \\
13.2.3 & La solución & Postulación en línea, puntaje con inteligencia artificial y auditoría solo de finalistas. & 0:50 \\
13.2.4 & Cómo funciona & El recorrido del postulante, de la DUBSS y del Directivo. & 0:50 \\
13.2.5 & Valor y avance & Equidad, eficiencia y trazabilidad; y qué está hecho. & 0:50 \\
13.2.6 & Equipo y cierre & Quiénes somos y qué pedimos. & 0:30 \\
\midrule
 & Total & & 4:00 \\
\bottomrule
\end{tabular}
\begin{tablenotes}
\item \textit{Nota.} Elaboración propia. Los tiempos son una guía y se ajustan al límite que fije el docente o el jurado.
\end{tablenotes}
\end{threeparttable}
\end{table}

\subsubsection{13.2.1. Gancho}

\begin{quote}
Imaginen que son uno de los 12\,000 estudiantes que entregaron su carpeta con todos sus documentos y que, por falta de tiempo, nadie la abrió jamás. ¿Es justo que una beca se decida por el orden en que llegó un papel?
\end{quote}

\subsubsection{13.2.2. El problema}

\begin{quote}
En la UAGRM la Beca Comedor ayuda a estudiantes que no pueden costear su almuerzo. Cada gestión llegan alrededor de 20\,000 carpetas físicas a la Dirección Universitaria de Bienestar Social y Salud, la DUBSS. Su personal logra revisar unas 8\,000, el 40\,\%. El otro 60\,\% nunca se abre. Y las becas disponibles son cerca de 900, una cifra que cambia con el presupuesto.

El resultado es un proceso lento, con riesgo de extravío de documentos sensibles, con sesgos y sin transparencia: un estudiante con gran necesidad puede quedar fuera simplemente porque su carpeta quedó al fondo de la pila.
\end{quote}

\subsubsection{13.2.3. La solución}

\begin{quote}
Proponemos el Sistema de Información Web y Móvil para la Postulación y Preselección Automatizada de la Beca Comedor, con reportes de Business Intelligence. La idea es simple: que la tecnología haga el trabajo repetitivo y las personas se dediquen a decidir.

La postulación pasa a ser cien por ciento en línea. Un módulo de inteligencia artificial calcula un puntaje de vulnerabilidad para todas las solicitudes, con datos como ingresos familiares, dependientes, procedencia y patrimonio. Con ese puntaje se ordenan los expedientes, y la DUBSS pide los documentos en papel y audita solo a los finalistas mejor puntuados: unas 2\,700 carpetas en lugar de 20\,000.
\end{quote}

\subsubsection{13.2.4. Cómo funciona}

\begin{quote}
El estudiante se registra, llena el formulario y adjunta sus documentos fotografiándolos con el celular. Puede seguir el estado de su postulación en cualquier momento y recibe avisos en el teléfono y por correo.

La DUBSS ve la lista priorizada, notifica a los finalistas y audita cada expediente. Al puntaje de la inteligencia artificial, que llega hasta 70 puntos, se suma el rendimiento académico que registra la persona auditora, hasta 30. Con el total, que llega a 100, decide a quién adjudicar la beca, publica los resultados y genera el acta.

Una vez que hay becados, el comedor registra la asistencia diaria desde el sistema, y el Directivo consulta un tablero con indicadores. La inteligencia artificial no decide sola: propone un orden, y las personas auditan y firman cada adjudicación.
\end{quote}

\subsubsection{13.2.5. Valor y avance}

\begin{quote}
Esto aporta tres cosas. Equidad: se evalúa el cien por ciento de las solicitudes, no el cuarenta. Eficiencia: se reciben, guardan y revisan unas 2\,700 carpetas físicas en lugar de 20\,000. Y trazabilidad: cada puntaje, auditoría y adjudicación queda registrado en una bitácora que se puede revisar.

El proyecto ya tiene sus requisitos, su análisis y su diseño documentados, y se desarrolla entre el 24 de agosto y el 30 de noviembre de 2026. Hoy están implementados los 26 casos de uso: la postulación, la preselección con inteligencia artificial, la auditoría, la adjudicación, la publicación de resultados, el acta, el control de asistencia y el tablero de indicadores. El sistema está en línea, y su API se probó caso de uso por caso de uso, sin fallos, con más de 3\,500 pruebas automatizadas escritas.
\end{quote}

% [VERIFICAR]: la última frase de 13.2.5 refleja el estado al 19-nov-2026 (26 CU, pruebas del Capítulo~\hyperref[cap:pruebas]{9}); si cambia antes de presentar, actualizarla.

\subsubsection{13.2.6. Equipo y cierre}

\begin{quote}
Somos NEXWARE, un equipo de dos estudiantes de la FICCT: Carlos Leonel Barrios Mayorga y Mauro Boris Gutierrez Primintela. Nuestro lema es «Si no existe, lo programamos».

Queremos pedirles dos cosas: validar con la DUBSS el puntaje mínimo y el número de cupos, y su respaldo para llevar este sistema a producción. Porque cada carpeta merece ser abierta.
\end{quote}

\subsection{13.3. Preguntas Frecuentes}

La Tabla~\ref{tab:pitch_preguntas} reúne las preguntas más probables y la respuesta corta que debe darse en cada caso.

\begin{longtable}{>{\raggedright\arraybackslash}p{4.8cm} >{\raggedright\arraybackslash}p{8.0cm}}
\caption{Preguntas previstas y respuestas breves}\label{tab:pitch_preguntas}\\
\toprule
Pregunta & Respuesta breve \\
\midrule
\endfirsthead
\toprule
Pregunta & Respuesta breve \\
\midrule
\endhead
\bottomrule
\endfoot
¿La inteligencia artificial decide quién recibe la beca? & No. Solo calcula un puntaje de vulnerabilidad que ordena los expedientes. La adjudicación la decide la DUBSS tras auditar el expediente físico del finalista. \\
¿Qué pasa si el puntaje de la inteligencia artificial se equivoca? & El puntaje se combina con el rendimiento académico y la auditoría de la DUBSS. Además, se registra con su fecha de cálculo y los parámetros del modelo pueden ajustarse y reprocesarse (CU-26). \\
¿Qué pasa con quien no es finalista? & Su solicitud fue evaluada con el mismo criterio que las demás y puede consultar su estado y resultado en el sistema. La apelación o reclamo está pendiente de consulta con la DUBSS. \\
¿Se elimina el papel? & No del todo. La documentación física se pide solo a los finalistas: unas 2\,700 carpetas en lugar de 20\,000. \\
¿Y la asistencia? & El sistema no tiene biometría propia: registra el resultado de la verificación del comedor y genera planillas imprimibles para las auditorías. \\
¿Cuál es el puntaje mínimo y cuántos cupos hay? & Son datos institucionales que se definen con la DUBSS en cada gestión; el sistema los toma como parámetros. \\
¿Qué tan seguros están los datos? & El acceso usa tokens y las contraseñas se guardan como hash, la comunicación es por HTTPS y los archivos se guardan en un almacenamiento propio. Las acciones relevantes quedan en la bitácora. \\
¿Puede usarse en otras instituciones? & Sí. El modelo de negocio previsto es ofrecer el sistema, adaptado, a otras instituciones con muchos postulantes y pocos cupos (Capítulo~\hyperref[cap:mercado]{8}). \\
\end{longtable}
\noindent\textit{Nota.} Elaboración propia a partir de los requisitos y del alcance del sistema.

\subsection{13.4. Recomendaciones para la Presentación}

\begin{itemize}
    \item Abrir con el gancho y con la cifra de 12\,000 carpetas sin abrir; es el dato que más se recuerda.
    \item Repetir tres cifras en toda la presentación: 20\,000 solicitudes, 900 cupos y 2\,700 carpetas físicas para finalistas.
    \item Mostrar el sistema en funcionamiento (recorrido del postulante y de la DUBSS) en lugar de describir pantallas (Capítulo~\hyperref[cap:producto]{12}).
    \item No prometer lo que aún no está implementado; presentarlo como parte del plan de desarrollo.
    \item Dividir el guion entre los dos integrantes: uno expone el problema y la solución, el otro el funcionamiento y el cierre.
\end{itemize}
